% Бібліографія гайду по геометрії. Записи проєкту взято з manual/bib/bibliography.tex і guide/bib/refs_project.tex.
\begin{thebibliography}{99}
\footnotesize\raggedright
\addcontentsline{toc}{chapter}{Бібліографія}

\bibitem{Manual} Фізика плазми всередині токамака~: навчальний посібник (методичка)~/ \texttt{manual/main.pdf}. -- 2026.
  Розд.~3 (рівновага), 4 (геометричні ефекти), 5 (анатомія токамака: JET).

\bibitem{Guide} Путівник гіпотезами проєкту~/ \texttt{guide/main.pdf}. -- 2026.

\bibitem{RegEst} Реєстр встановленого (E1--E35, R1--R9)~/ \texttt{Our try/00-decisions/ESTABLISHED.md}. -- Ред. 2026-09-29.

\bibitem{RegViz} Реєстр візуалізації (VE, VR, VC, VQ)~/ \texttt{tokamak-3d-viz/findings.md}. -- Ред. 2026-09-29.

\bibitem{ProjViz} tokamak-3d-viz~: візуалізація рівноваги DIII-D \#203702 (кадр 75, $t=1760$~мс),
  трасування силових ліній, острови, перерізи Пуанкаре~: програмний код (MIT)~/
  тека \texttt{tokamak-3d-viz}; опис конвенцій --- \texttt{docs/PHYSICS.md}. -- 2026. Дані рівноваги: Fusion Equilibrium Challenge
  (Sophelio / General Atomics), CC BY 4.0.

\bibitem{ProjJET} tokamak-3d-viz~: 3D-реконструкція JET (проєкт 1975 і стан 1983) --- база розмірів
  \texttt{machines/jet1975}, геометричне ядро, рівновага Соловйова, гофрування~: програмний код~/
  тека \texttt{tokamak-3d-viz}; опис --- \texttt{docs/JET-EQUILIBRIUM.md}. -- 2026.

\bibitem{FEC2026} Nakkina et al. The Fusion Equilibrium Challenge: Inferring Magnetic Geometry
  Without Magnetic Diagnostics~// arXiv:2609.01750. -- 2026. Набір даних
  \texttt{Sophelio/fusion-equilibrium-challenge}, CC BY 4.0.

\bibitem{JET1975} The JET Project~: Design Proposal for the Joint European Torus~: EUR 5516e (EUR-JET-R5). --
  Luxembourg~: Commission of the European Communities, 1976. -- 678~p. -- © CECA, CEE, CEEA.

\bibitem{Cerfon2010} Cerfon~A.\,J., Freidberg~J.\,P. «One size fits all» analytic solutions to the Grad--Shafranov
  equation~// Physics of Plasmas. -- 2010. -- Vol.~17, No.~3. -- 032502. -- DOI: 10.1063/1.3328818.
  (Поза корпусом; використано лише як метод у коді проєкту.)

\bibitem{Miller1998} Miller~R.\,L., Chu~M.\,S., Greene~J.\,M., Lin-Liu~Y.\,R., Waltz~R.\,E. Noncircular, finite aspect
  ratio, local equilibrium model~// Physics of Plasmas. -- 1998. -- Vol.~5, No.~4. -- P.~973--978.
  (Поза корпусом; джерело параметризації D-перерізу.)

\end{thebibliography}
